\section{Independent directions at depth four}
\label{oq:section}

The most recent ordered-Hurwitz report asks for an explicit description of
the first independent-slope case beyond depth three
\cite[Questions 1--2]{RepoOrdered}. We give such a description. Three
additional coordinates suffice beyond its previously defined ordering
function $\eta(a)$. All three have absolutely convergent representatives.
The resulting formula includes the complete principal part, every
transverse straight direction, exact shift recurrences, and a finite
reconstruction of the coordinates from three positive directions.
The construction is a spanning normal form; it does not assert arithmetic
independence of its coordinates.

\subsection{Conventions and the holomorphic regular germs}

Fix $a>0$ and use real logarithms of $n+a$. Set
\[
 Z_d(s_1,\ldots,s_d;a)
 =\sum_{n_1>\cdots>n_d\ge0}\prod_{j=1}^d(n_j+a)^{-s_j},
 \qquad Z_0=1.
\]
The initial domain is $\Re(s_1+\cdots+s_j)>j$ for every $j$.
Write $u_j=s_j-1$ and $P_k=u_1+\cdots+u_k$.
The Laurent convention is
\begin{equation}\label{oq:stieltjes}
 \zeta(1+u,a)=\frac1u+
 \sum_{j\ge0}\frac{(-1)^j\gamma_j(a)}{j!}u^j.
\end{equation}
Throughout this section,
\[
 g_j=\gamma_j(a),\quad g=g_0=-\psi(a),\quad
 z_k^{[j]}=\left.\partial_s^j\zeta(s,a)\right|_{s=k},
 \quad z_k=z_k^{[0]}.
\]
Thus a prime on an ordinary zeta value denotes its order derivative,
whereas $\eta'(a)$ below denotes its shift derivative.

For completeness, the regular germs can be constructed directly.
Define
\begin{align}
 D(s,x)&=\zeta(s,x+1)-\frac{x^{1-s}}{s-1},\label{oq:D}\\
 E_d(\mathbf u;a)&=
 \sum_{n_2>\cdots>n_d\ge0}
 D(1+u_1,n_2+a)\prod_{j=2}^d(n_j+a)^{-1-u_j}.
 \label{oq:E}
\end{align}
Let $R_0=1$, $R_1(u;a)=\zeta(1+u,a)-1/u$, and recursively put
\begin{align}\label{oq:Rrec}
 R_d(\mathbf u;a)=E_d(\mathbf u;a)+
 \frac{R_{d-1}(u_1+u_2,u_3,\ldots;a)
       -R_{d-1}(u_2,u_3,\ldots;a)}{u_1}.
\end{align}
The divided difference takes its removable value when $u_1=0$.

\begin{lemma}[Ordered decomposition and normalization]\label{oq:regular}
The germs $R_d$ are holomorphic on $\sum_j|u_j|<1/2$, and
\begin{equation}\label{oq:decomp}
 Z_d(1+\mathbf u;a)=
 \sum_{k=0}^d\frac{R_{d-k}(u_{k+1},\ldots,u_d;a)}{P_1\cdots P_k}.
\end{equation}
Furthermore, if $B_d(a)=R_d(0,\ldots,0;a)$, then
\begin{equation}\label{oq:gamma}
 \sum_{d\ge0}B_d(a)y^d=\frac{\Gamma(a)}{\Gamma(a+y)}.
\end{equation}
\end{lemma}
\begin{proof}
For $\Re s>0$, the tail defect has the integral
\[
 D(s,x)=\frac1{\Gamma(s)}\int_0^\infty
 t^{s-1}e^{-xt}\left(\frac1{e^t-1}-\frac1t\right)dt.
\]
The bracket is bounded at zero. This proves removability at $s=1$
and $D(s,x)=O(x^{-\Re s})$ on compact subsets. The summands in
\eqref{oq:E}, after summing the inner finite indices, have a
uniform bound of the form
$C(1+n_2)^{-2+\sum|u_j|}(1+\log(1+n_2))^{d-2}$.
Normal convergence and the integral expression for a divided difference
prove holomorphy. Summing first over $n_1>n_2$ gives
\[
 Z_d(1+\mathbf u;a)=\frac1{u_1}
 Z_{d-1}(1+u_1+u_2,1+u_3,\ldots;a)+E_d(\mathbf u;a),
\]
so induction proves \eqref{oq:decomp}. On the common diagonal,
the elementary-symmetric product for $Z_d$ implies, coefficientwise
as a formal power series in $y$,
\begin{align}\label{oq:diaggenerator}
 \sum_{d\ge0}R_d(t,\ldots,t;a)y^d
 =\exp\left\{yR_1(t;a)+
 \sum_{m\ge2}\frac{(-1)^{m-1}}m\zeta(m+mt,a)y^m\right\}.
\end{align}
At $t=0$, the logarithmic Taylor expansion of $\Gamma(a+y)$
gives \eqref{oq:gamma}.
\end{proof}

This construction and the diagonal generating function are antecedents,
not new claims of this article. For $a=1$, the ordered Laurent framework
is due to Saha \cite{Saha}; general integer-point Laurent expansions are
studied in \cite{MOW}. The common-shift construction above is the one
in \cite{RepoOrdered}. In particular,
\begin{align}
 B_2&=\frac{g^2-z_2}{2},&
 B_3&=\frac{g^3-3gz_2+2z_3}{6},\label{oq:B23}\\
 B_4&=\frac{g^4-6g^2z_2+3z_2^2+8gz_3-6z_4}{24}.
 \label{oq:B4}
\end{align}

\subsection{The small jet that controls the answer}

Define
\begin{align}
 \eta(a)&=[u]R_2(u,0;a),\label{oq:eta}\\
 \rho(a)&=[u^2]R_2(u,0;a)-[v^2]R_2(0,v;a),\label{oq:rho}\\
 \kappa(a)&=(\partial_1-2\partial_2+\partial_3)R_3(0,0,0;a),
 \label{oq:kappa}\\
 T(a)&=(\partial_u-\partial_v)Z_2(2+u,1+v;a)\big|_{u=v=0}.
 \label{oq:T}
\end{align}
The factorial convention in \eqref{oq:rho} matters: $\rho$ is a
difference of Taylor coefficients, not a difference of second derivatives.
The last coordinate has the elementary convergent series
\begin{equation}\label{oq:Tseries}
 T(a)=\sum_{n>m\ge0}
 \frac{\log(m+a)-\log(n+a)}{(n+a)^2(m+a)}.
\end{equation}
Introduce the ordinary-data combinations
\begin{align}
 D_1&=-gg_1-z_2^{[1]},\label{oq:D1}\\
 E&=\frac{gg_2-z_2^{[2]}}2,
 &F&=\frac{g_1^2-z_2^{[2]}}2,\label{oq:EF}\\
 S&=-g_1B_2-gz_2^{[1]}+z_3^{[1]},\label{oq:S}\\
 V&=g\eta+\frac{g_1(g^2+z_2)}2+gz_2^{[1]}-T.
 \label{oq:V}
\end{align}
Only $V$ retains a nonordinary coordinate; the letters $D_1,E,F,S$
denote explicitly displayed finite expressions.

\begin{theorem}[Complete required jet]\label{oq:jet}
The quadratic Taylor polynomial of $R_2$ and the linear Taylor polynomial
of $R_3$ are
\begin{align}\label{oq:R2jet}
 R_2(u,v;a)={}&B_2+\eta u+(D_1-\eta)v
 +\frac{E+\rho}{2}u^2+Fuv+\frac{E-\rho}{2}v^2
 +O((|u|+|v|)^3),\\
 R_3(u,v,w;a)={}&B_3+
 \frac S3(u+v+w)+\frac\kappa6(u-2v+w)
 +\frac V2(u-w)+O((|u|+|v|+|w|)^2).
 \label{oq:R3jet}
\end{align}
\end{theorem}
\begin{proof}
The regularized depth-two stuffle identity is
\begin{equation}\label{oq:stuffle2}
 R_2(u,v)+R_2(v,u)=R_1(u)R_1(v)-\zeta(2+u+v,a).
\end{equation}
It follows directly by substituting \eqref{oq:decomp} into the ordinary
stuffle identity; the rational polar terms cancel. Its coefficients of
$u$, $u^2$ and $uv$, together with the definition of $\rho$, prove
\eqref{oq:R2jet}.

For the next depth put
\[
 C(u,v)=Z_2(2+u,1+v;a),\qquad
 L(u,v)=Z_2(1+u,2+v;a)-\frac{\zeta(2+v,a)}u.
\]
The first function converges normally near zero. The second is holomorphic
there because the two-factor stuffle law gives the exact expression
\begin{equation}\label{oq:L}
 L(u,v)=R_1(u)\zeta(2+v,a)-C(v,u)-\zeta(3+u+v,a).
\end{equation}
Applying the three-factor stuffle identity and cancelling the explicit
prefix poles gives
\begin{align}\label{oq:stuffle3}
 R_1(u)R_2(v,w)={}&R_3(u,v,w)+R_3(v,u,w)+R_3(v,w,u)
 +C(u+v,w)+L(v,u+w).
\end{align}
For clarity, the needed cancellation uses
$R_2(u,w)+R_2(w,u)=R_1(u)R_1(w)-\zeta(2+u+w,a)$.
Thus \eqref{oq:stuffle3} follows from meromorphic identities and entails
no unproved assumption that arbitrary finite-part prescriptions preserve
products.

Write $X=\partial_1R_3(0)$, $Y=\partial_2R_3(0)$ and
$Z=\partial_3R_3(0)$. Differentiating \eqref{oq:stuffle3} with respect
to $u$ and using \eqref{oq:L} gives
\[
 X+Y+Z=-g_1B_2-gz_2^{[1]}+z_3^{[1]}=S.
\]
Differentiating with respect to $v$ instead gives
\[
 2X+Y=g\eta+g_1z_2+z_3^{[1]}-T.
\]
Subtracting the preceding equation proves $X-Z=V$.
Finally $X-2Y+Z=\kappa$ by definition. Solving these three linear
equations gives \eqref{oq:R3jet}.
\end{proof}

\subsection{Every transverse depth-four Laurent constant}

For $p,b,c,d\in\mathbb C$ suppose
\begin{equation}\label{oq:transverse}
 p(p+b)(p+b+c)(p+b+c+d)\ne0.
\end{equation}
Here $\operatorname{FP}_{\varepsilon=0}$ means the constant coefficient
of the meromorphic function after the specified substitution. It is not
an iterated residue and is not a value on a polar hyperplane.

\begin{theorem}[Depth-four directional normal form]\label{oq:main}
Under \eqref{oq:transverse},
\begin{align}\label{oq:FP4}
 &\operatorname{FP}_{\varepsilon=0}
 Z_4(1+p\varepsilon,1+b\varepsilon,1+c\varepsilon,1+d\varepsilon;a)
 \notag\\
 &\quad=B_4+
 \frac1p\left\{\frac{S(b+c+d)}3+
            \frac{\kappa(b-2c+d)}6+\frac{V(b-d)}2\right\}
 \notag\\
 &\qquad\quad+
 \frac{E(c^2+d^2)/2+Fcd+\rho(c^2-d^2)/2}{p(p+b)}
 -\frac{g_3d^3}{6p(p+b)(p+b+c)}.
\end{align}
The complete principal part is
\begin{align}\label{oq:principal}
 &\frac{1}{p(p+b)(p+b+c)(p+b+c+d)\varepsilon^4}
 +\frac{g}{p(p+b)(p+b+c)\varepsilon^3}
 \notag\\
 &\quad+\frac1{\varepsilon^2}
 \left\{\frac{B_2}{p(p+b)}-
 \frac{dg_1}{p(p+b)(p+b+c)}\right\}
 \notag\\
 &\quad+\frac1\varepsilon
 \left\{\frac{B_3}p+
 \frac{\eta(c-d)+D_1d}{p(p+b)}+
 \frac{g_2d^2}{2p(p+b)(p+b+c)}\right\}.
\end{align}
\end{theorem}
\begin{proof}
At depth four, \eqref{oq:decomp} has five terms:
\[
 \frac1{P_1P_2P_3P_4}+
 \frac{R_1(u_4)}{P_1P_2P_3}+
 \frac{R_2(u_3,u_4)}{P_1P_2}+
 \frac{R_3(u_2,u_3,u_4)}{P_1}+R_4(\mathbf u).
\]
The constant term requires the cubic coefficient of $R_1$, the quadratic
coefficient of $R_2$, the linear coefficient of $R_3$, and $R_4(0)=B_4$.
Equations \eqref{oq:stieltjes}, \eqref{oq:R2jet} and \eqref{oq:R3jet}
give exactly \eqref{oq:FP4}. The lower powers give
\eqref{oq:principal}. Holomorphy of the remainders makes the coefficient
calculation legitimate, including directions with some zero or negative
slopes, provided the prefix sums remain nonzero.
\end{proof}

The answer depends on $p$ and the three trailing slopes, but it needs
no new fourth-depth regular constant. The new information lies in
lower-depth derivatives, exactly as the ordered pole geometry predicts.

\begin{corollary}[Exact cancellation locus in this normal form]
\label{oq:locus}
Treat $\rho,\kappa,T,\eta$ as named coordinates, without assigning
unproved relations among their values. Under \eqref{oq:transverse},
all their coefficients in \eqref{oq:FP4} vanish if and only if
$b=c=d$. On this locus the answer reduces to ordinary Hurwitz jets.
\end{corollary}
\begin{proof}
The coefficient of $T$ is $-(b-d)/(2p)$ and the coefficient of
$\kappa$ is $(b-2c+d)/(6p)$. Their simultaneous vanishing forces
$b=d$ and then $c=b$. This also kills the $\rho$ and $\eta$
coefficients. Conversely, substitution proves cancellation.
\end{proof}

This corollary classifies cancellation in the proved coordinate expression.
It does not say that no other direction can simplify after a new functional
identity is discovered. In particular, it is not a statement of arithmetic
independence at an individual value of $a$.

For $b=c=d=q$, the formula becomes
\begin{equation}\label{oq:diagonalcheck}
 B_4+\frac qpS+
 \frac{q^2}{p(p+q)}\left\{\frac{g_1^2+gg_2}{2}-z_2^{[2]}\right\}
 -\frac{q^3g_3}{6p(p+q)(p+2q)},
\end{equation}
which recovers the already proved distinguished-slope formula of
\cite{RepoOrdered}. This is a consistency specialization, not an additional
new evaluation.

\subsection{Three positive rays recover all three coordinates}

Define the ordinary part of \eqref{oq:FP4} by
\begin{align}\label{oq:ordinary}
 \mathcal O(p,b,c,d;a)=B_4+
 \frac{S(b+c+d)}{3p}+
 \frac{E(c^2+d^2)/2+Fcd}{p(p+b)}
 -\frac{g_3d^3}{6p(p+b)(p+b+c)}.
\end{align}
Let $\Delta_A,\Delta_B,\Delta_C$ be the directional finite parts minus
$\mathcal O$ at the respective rays
\[
 A=(1,1,2,1),\qquad B=(1,1,3,1),\qquad C=(1,2,1,1).
\]
All three lie in the positive convergence directions.

\begin{corollary}[Finite reconstruction]\label{oq:reconstruction}
One has
\begin{align}
 \rho&=2\Delta_B-4\Delta_A,\label{oq:recoverrho}\\
 \kappa&=\frac92\Delta_B-12\Delta_A,\label{oq:recoverkappa}\\
 V&=2\Delta_C-\frac32\Delta_B+4\Delta_A.\label{oq:recoverV}
\end{align}
Consequently $T$ is recovered from \eqref{oq:V} once $\eta$ is known.
\end{corollary}
\begin{proof}
The three equations supplied by \eqref{oq:FP4} are
\[
 \begin{pmatrix}\Delta_A\\\Delta_B\\\Delta_C\end{pmatrix}
 =\begin{pmatrix}-1/3&3/4&0\\-2/3&2&0\\1/6&0&1/2\end{pmatrix}
 \begin{pmatrix}\kappa\\\rho\\V\end{pmatrix}.
\]
The determinant is $-1/12$. Inverting this rational matrix gives the
displayed formulas.
\end{proof}

Thus the appearance of three coordinates is observable by a finite exact
linear calculation. The invertible matrix proves equivalence of two
spanning descriptions; it does not prove that the three numerical values
are independent over any field.

\subsection{Absolutely convergent representatives}

Put
\[
 Q_j(x)=\gamma_j(x+1)+\frac{\log^{j+1}x}{j+1},\qquad
 h_n(a)=\sum_{m=0}^{n-1}\frac1{m+a},\qquad
 \ell_n(a)=\sum_{m=0}^{n-1}\frac{\log(m+a)}{m+a}.
\]
Here empty sums vanish. The familiar identities
\[
 h_n(a)=\psi(n+a)-\psi(a),\qquad
 \ell_n(a)=\gamma_1(a)-\gamma_1(n+a)
\]
make the representatives numerically usable with one outer summation.

\begin{theorem}[Convergent coordinate formulas]\label{oq:series}
With $x=n+a$ in every sum,
\begin{align}
 \eta(a)&=\frac{g_2}{2}-\sum_{n\ge0}\frac{Q_1(x)}x,
 \label{oq:etaseries}\\
 \rho(a)&=\frac{g_3}{3}+
 \frac12\sum_{n\ge0}\frac{Q_2(x)-Q_0(x)\log^2x}{x},
 \label{oq:rhoseries}\\
 \kappa(a)&=\sum_{n\ge0}
 \frac{-Q_1(x)h_n(a)+Q_0(x)(2\log x\,h_n(a)-\ell_n(a))}{x}
 -\frac32(E+\rho)+F.\label{oq:kappaseries}
\end{align}
All these series, as well as \eqref{oq:Tseries}, converge absolutely
and locally uniformly for $a$ in compact subintervals of $(0,\infty)$.
\end{theorem}
\begin{proof}
Taylor expansion in \eqref{oq:D} gives
$D(1+u,x)=\sum_{j\ge0}(-1)^jQ_j(x)u^j/j!$.
If $r^{(d)}_{m_1,\ldots,m_d}=[\mathbf u^{\mathbf m}]R_d$, the
divided difference in \eqref{oq:Rrec} yields
\begin{equation}\label{oq:coeffrec}
 r^{(d)}_{m_1,\ldots,m_d}
 =[\mathbf u^{\mathbf m}]E_d+
 \binom{m_1+m_2+1}{m_1+1}
 r^{(d-1)}_{m_1+m_2+1,m_3,\ldots,m_d}.
\end{equation}
For $(m_1,m_2)=(1,0),(2,0),(0,2)$ this gives the first two formulas.
For $(1,0,0),(0,1,0),(0,0,1)$, their combination with coefficients
$1,-2,1$ gives the sum in \eqref{oq:kappaseries} and the correction
$-3[u^2]R_2+[uv]R_2=-3(E+\rho)/2+F$.
The estimates used for \eqref{oq:E}, followed by Cauchy's formula on
a slightly larger compact spectral polydisc, justify every coefficient
extraction and give absolute convergence. Alternatively,
$Q_j(x)=O((1+\log x)^j/x)$ suffices, since
$h_n=O(1+\log n)$ and $\ell_n=O(1+\log^2n)$.
\end{proof}

These formulas define the coordinates independently of a numerical
Laurent fit. They are useful even if no further ordinary-zeta reduction
exists, and they keep the unevaluated part confined to explicitly
convergent harmonic--Stieltjes sums.

\subsection{Exact shift and differential compatibility}

\begin{proposition}[Triangular shift laws]\label{oq:shift}
Let a superscript $+$ mean evaluation at $a+1$, and let $L=\log a$.
Then
\begin{align}
 \eta(a)-\eta(a+1)&=-\frac{g_1^+}{a},\label{oq:shifteta}\\
 \rho(a)-\rho(a+1)&=\frac{g_2^+-g^+L^2}{2a},\label{oq:shiftrho}\\
 T(a)-T(a+1)&=\frac{z_2^{[1],+}+Lz_2^+}{a},\label{oq:shiftT}\\
 \kappa(a)-\kappa(a+1)&=
 \frac{3\eta^+-2D_1^+-LB_2^+}{a}.\label{oq:shiftkappa}
\end{align}
\end{proposition}
\begin{proof}
The terms with last index zero give the exact ordered shift identity
\[
 Z_d(\mathbf s;a)-Z_d(\mathbf s;a+1)
 =a^{-s_d}Z_{d-1}(s_1,\ldots,s_{d-1};a+1).
\]
Substitution into \eqref{oq:decomp} and induction on depth show that
\begin{equation}\label{oq:Rshift}
 R_d(\mathbf u;a)-R_d(\mathbf u;a+1)
 =a^{-1-u_d}R_{d-1}(u_1,\ldots,u_{d-1};a+1).
\end{equation}
The first, second and fourth assertions are the indicated Taylor
coefficients of \eqref{oq:Rshift}. For the third, use
$C(u,v;a)-C(u,v;a+1)=a^{-1-v}\zeta(2+u,a+1)$ and apply
$\partial_u-\partial_v$.
\end{proof}

\begin{proposition}[A convergent derivative identity]\label{oq:derivative}
Let $C_0(a)=Z_2(2,1;a)$. Then
\begin{equation}\label{oq:etaT}
 T(a)=-\eta'(a)-C_0(a)+g_1z_2+z_3^{[1]}.
\end{equation}
At $a=1$, the classical Euler evaluation $Z_2(2,1;1)=\zeta(3)$ gives
\[
 T(1)=-\eta'(1)-\zeta(3)+\gamma_1\zeta(2)+\zeta'(3).
\]
\end{proposition}
\begin{proof}
Differentiating the ordered defining sum in its initial domain, and
then continuing, gives
\[
 \partial_aR_2(u,v;a)=-(1+u)C(u,v)-(1+v)L(u,v).
\]
The pole removed from $L$ cancels exactly with
$\partial_aR_1(v)/u=-(1+v)\zeta(2+v,a)/u$.
Apply $\partial_u$ at zero and substitute
$L_u(0,0)=-g_1z_2-z_3^{[1]}-C_v(0,0)$ from \eqref{oq:L}.
The result is \eqref{oq:etaT}. Normal convergence justifies the
derivatives of $C$; the regular germs are holomorphic in $a$ for
$\Re a>0$ on the principal branch.
\end{proof}

Equations \eqref{oq:shifteta}--\eqref{oq:etaT} connect the new directional
coordinates to generalized harmonic sums, generalized Stieltjes constants,
and parameter derivatives. They also supply independent consistency checks
for implementations. A complete reduction of $\rho$ or $\kappa$ to ordinary
single-zeta data is not proved here.
